\subsection{Enzyme representation}
\label{sec:enzyme-representation}
\label{sec:enzyme-encoder}
\label{sec:paper-v4-encoders}
Catalytic function depends on both sequence context and functional residues.
We capture these signals using frozen
ProtT5 residue embeddings $\mathbf{H}_e\in\mathbb R^{L\times1024}$
\citep{elnaggar2022prottrans} and scores from a frozen SLEEC model
\citep{gong2026sleec}, complemented by $K$ learned residue summaries.
The global ProtT5 feature is
$\mathbf{x}_e^{P}=L^{-1}\sum_{i\in\mathcal V_e}\mathbf{h}_i$,
where $\mathcal V_e$ indexes the $L$ valid residues.

\myparagraph{SLEEC-guided pooling}
We suppress $e$ on residue vectors and attention weights.
SLEEC assigns each $\mathbf{h}_i\in\mathbb R^{1024}$
a logit $\ell_i$, centered and clipped as
$\tilde\ell_i=\operatorname{clip}(\ell_i-L^{-1}\sum_{j\in\mathcal V_e}\ell_j,-10,10)$.
The SLEEC feature $\mathbf{x}_e^{S}\in\mathbb R^{1024}$ summarizes enzyme $e$
by averaging its residue embeddings with normalized SLEEC weights:
\begin{equation}
 a_i^{\mathrm S}=(1-\eta)
 \frac{\exp(\tilde\ell_i)}{\sum_{j\in\mathcal V_e}\exp(\tilde\ell_j)}
 +\frac{\eta}{L},\qquad
 \mathbf{x}_e^{S}=\sum_{i\in\mathcal V_e}a_i^{\mathrm S}\mathbf{h}_i,
 \label{eq:sleec-attention}
\end{equation}
Here $\eta=0.05$ retains a uniform contribution from every valid residue.

\myparagraph{Learned residue views}
We learn queries $\mathbf{q}_k\in\mathbb R^{256}$, $k=1,\ldots,K$, shared
across enzymes. A shared residue adapter and separate key and value projections
produce normalized keys $\mathbf{k}_i$ and unnormalized values $\mathbf{v}_i$
in $\mathbb R^{256}$. Write $\vunit{\mathbf{x}}=\mathbf{x}/\max(\|\mathbf{x}\|_2,\epsilon)$
for numerical normalization, with $\epsilon=10^{-6}$ in this encoder.
Query $\hat{\mathbf{q}}_k=\vunit{\mathbf{q}_k}$ defines
\begin{equation}
 p_{ki}=\frac{\exp(\gamma\hat{\mathbf{q}}_k^\top\mathbf{k}_i)}
 {\sum_{j\in\mathcal V_e}\exp(\gamma\hat{\mathbf{q}}_k^\top\mathbf{k}_j)},
 \qquad
 \mathbf{x}_e^{k}=\sum_{i\in\mathcal V_e}
 \left[(1-\eta)p_{ki}+\frac{\eta}{L}\right]\mathbf{v}_i,
 \label{eq:paper-learned-views}
\end{equation}
Thus $p_{ki}$ is a scalar weight over residues for query $k$; the learned scale
$\gamma=10\vsigmoid(\rho_\gamma)\in(0,10)$ controls concentration. Each view
combines 95\% attention-weighted and 5\% uniform pooling of the values.
Phase 1 trains the queries, adapter, and projections; frozen views can be
cached per enzyme for reuse across candidate reactions. We retain
$K=4$ as a compromise across screening and retrieval: the single-seed
$K\in\{1,2,4,8\}$ ablation finds no uniformly best count
(Appendix~\ref{app:v4-k-ablation}).

\myparagraph{Feature-wise gated fusion}
We combine the extracted enzyme features
$\{\mathbf{x}_e^{P},\mathbf{x}_e^{S}\}\cup\{\mathbf{x}_e^{k}\}_{k=1}^{K}$
by projecting each feature into a common 256-dimensional space using an
independent MLP followed by $\ell_2$ normalization. The projected features are
$\mathbf{y}_e^{v}=\vunit{F_v(\mathbf{x}_e^{v})}$, with view labels
$v\in\mathcal I=\{P,S,1,\ldots,K\}$.
A gating network maps their concatenation to feature-specific weights
$\omega_{e,vd}$, normalized across views for each coordinate $d$.
The fused representation is
\begin{equation}
 \bar y_{e,d}=\sum_{v\in\mathcal I}\omega_{e,vd}y_{e,d}^{v},
 \qquad d=1,\ldots,256.
 \label{eq:paper-protein}
\end{equation}
Independent projections of $\mathbf{x}_e^{P}$ and $\bar{\mathbf{y}}_e$ yield
normalized 512-dimensional vectors $\mathbf{g}_e$ and $\mathbf{f}_e$.
The enzyme embedding is $\mathbf{b}_e=\vunit{\mathbf{g}_e+s\mathbf{f}_e}$,
where the learned scalar $s\in(0,0.5)$ controls the fused contribution
(Appendix~\ref{sec:v4-protein}).
